\section{The Canonical Split of the Support Register}

The support register
\[
\mathbf h=(h_1,\ldots,h_6)
\]
is a vector in \(\mathbb R^6\). This section separates it into a common
component shared by all six coordinates and a complementary zero-mean
component.

The decomposition is standard linear algebra. Its Euclidean meaning is
proved in Section~4.

\subsection{Uniform and zero-mean subspaces}

Let
\[
\mathbf 1=(1,1,1,1,1,1)
\]
and define the uniform subspace
\[
U=\operatorname{span}\{\mathbf 1\}.
\]

A register is uniform precisely when all six coordinates are equal.

Define also
\[
H=
\left\{
\mathbf r\in\mathbb R^6:
\sum_{i=1}^{6}r_i=0
\right\}.
\]

The subspace \(H\) is the kernel of the nonzero linear functional
\[
\mathbf x\longmapsto\sum_{i=1}^{6}x_i,
\]
and therefore
\[
\dim H=5.
\]

\subsection{Mean support and residual}

For any
\[
\mathbf h=(h_1,\ldots,h_6)\in\mathbb R^6,
\]
define its mean support by
\[
\bar h
=
\frac{1}{6}\sum_{i=1}^{6}h_i.
\]

The corresponding uniform component is
\[
\mathbf h_U
=
\bar h\,\mathbf 1.
\]

Define the residual by
\[
\mathbf r
=
\mathbf h-\bar h\,\mathbf 1.
\]

Its coordinate sum vanishes:
\[
\sum_{i=1}^{6}r_i
=
\sum_{i=1}^{6}h_i-6\bar h
=
0.
\]

Hence
\[
\mathbf r\in H.
\]

\subsection{Canonical decomposition}

\begin{theorem}[Canonical decomposition]
Every vector \(\mathbf h\in\mathbb R^6\) admits a unique decomposition
\[
\mathbf h
=
\bar h\,\mathbf 1+\mathbf r,
\]
where
\[
\bar h\,\mathbf 1\in U
\qquad\text{and}\qquad
\mathbf r\in H.
\]
\end{theorem}

\begin{proof}
The definitions above give existence.

For uniqueness, suppose
\[
\mathbf h
=
\mathbf u_1+\mathbf r_1
=
\mathbf u_2+\mathbf r_2,
\]
with
\[
\mathbf u_1,\mathbf u_2\in U
\qquad\text{and}\qquad
\mathbf r_1,\mathbf r_2\in H.
\]

Then
\[
\mathbf u_1-\mathbf u_2
=
\mathbf r_2-\mathbf r_1
\in U\cap H.
\]

If \(a\mathbf 1\in U\cap H\), then
\[
0=\sum_{i=1}^{6}a=6a,
\]
so \(a=0\). Thus
\[
U\cap H=\{\mathbf 0\},
\]
and both components are unique.
\end{proof}

Therefore
\[
\mathbb R^6=U\oplus H.
\]

\subsection{Orthogonality}

Equip \(\mathbb R^6\) with the standard inner product
\[
\langle\mathbf x,\mathbf y\rangle
=
\sum_{i=1}^{6}x_i y_i.
\]

For every \(\mathbf r\in H\),
\[
\langle\mathbf 1,\mathbf r\rangle
=
\sum_{i=1}^{6}r_i
=
0.
\]

Hence
\[
H=U^\perp,
\]
and the direct sum is orthogonal:
\[
\mathbb R^6
=
U\mathbin{\perp}H.
\]

In particular,
\[
\left\langle
\bar h\,\mathbf 1,
\mathbf r
\right\rangle
=
0.
\]

Thus
\[
\|\mathbf h\|^2
=
6\bar h^2+\|\mathbf r\|^2.
\]

\subsection{Projection operators}

Let \(J\) be the \(6\times6\) all-ones matrix. Define
\[
P_U=\frac{1}{6}J
\]
and
\[
P_H=I-\frac{1}{6}J.
\]

Then
\[
P_U\mathbf h
=
\bar h\,\mathbf 1
\]
and
\[
P_H\mathbf h
=
\mathbf r.
\]

Since
\[
J^2=6J,
\]
the operators satisfy
\[
P_U^2=P_U,
\qquad
P_H^2=P_H,
\]
\[
P_UP_H=P_HP_U=0,
\]
and
\[
P_U+P_H=I.
\]

They are therefore the complementary orthogonal projections onto \(U\)
and \(H\).

\subsection{Symmetry compatibility}

Every rotational symmetry of the regular dodecahedron acts on support space
by a permutation matrix \(\rho(g)\).

Coordinate permutations fix \(\mathbf 1\) and preserve coordinate sums.
Therefore
\[
\rho(g)U=U
\qquad\text{and}\qquad
\rho(g)H=H.
\]

The projection operators commute with the symmetry action:
\[
\rho(g)P_U=P_U\rho(g),
\qquad
\rho(g)P_H=P_H\rho(g).
\]

Consequently,
\[
P_U(\rho(g)\mathbf h)
=
\rho(g)(P_U\mathbf h)
\]
and
\[
P_H(\rho(g)\mathbf h)
=
\rho(g)(P_H\mathbf h).
\]

The split is therefore independent of the arbitrary labeling of the six
opposite-face pairs.

\subsection{Positive support registers}

The algebraic decomposition applies to every vector in \(\mathbb R^6\).
The geometric support realization requires
\[
\mathbf h\in\mathbb R_{>0}^6.
\]

If
\[
\mathbf h
=
\bar h\,\mathbf 1+\mathbf r
\]
is positive, then
\[
\bar h+r_i>0
\]
for every \(i\), and necessarily
\[
\bar h>0.
\]

The residual may contain positive and negative entries. These are deviations
from a positive mean support, not negative support distances.

\subsection{Interpretation boundary}

The canonical split establishes one uniform coordinate and five independent
zero-mean coordinates.

It does not by itself prove that the uniform component is Euclidean scale
or that the residual is geometric shape variation.

Section~4 proves those conclusions from the support realization.
